\begin{abstract}
\textbf{[Draft abstract on preliminary laptop-scale results; numbers will change.]}
The whole-brain leaky integrate-and-fire model of Shiu et al.\ predicts feeding and grooming
circuits from the FlyWire connectome of one \textit{Drosophila}. Earlier ports found that
its predictions transfer poorly to a second female connectome (BANC), but could not say
whether the cause is biological or technical. We ran the unchanged model on FlyWire, BANC
and the male CNS with a float32 implementation validated against the original Brian2 code
(rate correlations 0.9994--0.9997). Per cell type, BANC carries 0.45 times and the male CNS
1.76 times FlyWire's input synapses. With raw weights BANC lost sugar$\to$MN9 and
JO$\to$aBN1, and the male CNS reproduced them only with brain-wide activity. Cell-type
response profiles agreed weakly between connectomes (median Spearman $\rho$ 0.03--0.21)
and better between hemispheres of FlyWire (0.35) and the male CNS (0.52). Thinning FlyWire
to BANC's per-type input abolished sugar$\to$MN9, so reconstruction density alone explains
that failure. It did not explain the rest: global rescaling rescued the thinned FlyWire but
not BANC, and thinned FlyWire stayed closer to FlyWire ($\rho\approx0.5$) than BANC was.
Completeness therefore explains part, but not all, of the cross-connectome gap.
\todo{update after the full grid}
\end{abstract}
